% =====================================================================
\chapter{Tercer instrumento: el VIX argentino y el panel de acciones}
\chaptermark{El VIX argentino y el panel de acciones}
\label{cap:vix}

El VIX de Chicago es el termómetro más conocido del miedo en los mercados: sube cuando los
inversores esperan movimientos bruscos. Se calcula con precios de opciones, y en Argentina no
existe un mercado de opciones con la liquidez necesaria para construir algo así. Por eso se
construyó un \textit{VIX argentino} a partir de lo que sí está disponible todos los días: los
retornos del S\&P~Merval. Es el instrumento que mide cuánta tensión hay en el mercado, y la
base sobre la que se estiman los regímenes del capítulo~\ref{cap:hmm}. Este capítulo presenta
también el panel de acciones del lado bursátil: cómo se eligió, cómo se depuraron sus precios,
cuán volátil es cada acción y si las acciones forman familias de comportamiento.

% ---------------------------------------------------------------------
\section{El panel de acciones}
\label{sec:panel}

Los precios diarios (apertura, máximo, mínimo, cierre y volumen) se descargaron de Yahoo Finance
con \texttt{yfinance}. Se partió de 53 candidatas del Panel General y del Panel Líderes de BYMA
(las 20 del S\&P~Merval más 33 de sectores similares) y se exigieron dos condiciones en los tres
años de datos (731 ruedas): cobertura completa del calendario y al menos 90\% de días con volumen
real, porque un día sin operaciones no es un faltante sino una fila con el precio pegado. Pasaron
47, incluidas las 20 del índice; quedaron afuera Carboclor, Holcim, Polledo, Instituto Rosenbusch
y dos tickers sin datos (NII Holdings y TGLT, que hoy cotiza como GCDI). En la depuración se
quitaron las filas que Yahoo agrega en los feriados (se usó como calendario de BYMA la lista de
fechas del índice), se reintentaron los tickers con fallas intermitentes y se corrigieron dos
eventos corporativos mal ajustados: el \textit{split} de SEMI del 16/12/2024 y el dividendo en
acciones de AGRO de noviembre de 2023, que mostraba una caída falsa de 88\%.

% ---------------------------------------------------------------------
\section{Construcción del VIX argentino}

El VIX argentino es la \textit{volatilidad realizada} del índice: el desvío estándar de sus
retornos diarios en las últimas 20 ruedas, anualizado \citep{sakowski2023}:
\begin{equation}
\sigma_t = \sqrt{252}\cdot \operatorname{desv}\left(r_{t-19},\dots,r_t\right)\cdot 100,
\qquad r_t = \frac{P_t - P_{t-1}}{P_{t-1}},
\end{equation}
donde $P_t$ es el cierre del S\&P~Merval (\texttt{\^{}MERV} en Yahoo Finance) en la rueda $t$.
Veinte ruedas equivalen a un mes de operaciones, y el factor $\sqrt{252}$ es la convención
internacional de ruedas por año, que lleva el desvío diario a una escala anual comparable con
la de cualquier otro índice (BYMA tiene unas 244 ruedas por año por sus feriados; usar ese
número bajaría todos los valores un 1,5\%, sin cambiar ninguna comparación). Se calculó para la ventana del corpus de noticias: del 2 de enero
de 2025 al 30 de abril de 2026, 323 ruedas. El mismo cálculo, aplicado a cada acción, da su
volatilidad propia. La lista de fechas del índice se usó, además, como calendario de ruedas de
BYMA para todo el trabajo.

% ---------------------------------------------------------------------
\section{Cómo leer el gráfico}

La figura~\ref{fig:vix} tiene dos paneles que se leen juntos, y las líneas punteadas marcan los
tres hechos que dispararon la volatilidad del período: la salida del cepo y las dos elecciones. Arriba está el S\&P~Merval en
pesos: como los precios son nominales, su tendencia de fondo incluye la inflación y el tipo de
cambio, y lo que interesa son los saltos, no la pendiente. Abajo está el VIX argentino, con su
promedio del período (44,0\%) como referencia; el sombreado marca los días en que estuvo por
encima de ese promedio (136 de las 323 ruedas). Un VIX de 44\% quiere decir que, al ritmo del
último mes, el índice se movería en un año alrededor de un 44\% hacia arriba o hacia abajo. El
mínimo fue de 15,4\% el 30/04/2026 y el máximo, de 91,0\% el 21/11/2025.

\begin{figure}[H]
\centering
\includegraphics[width=0.88\linewidth]{vix_arg.png}
\caption{S\&P~Merval en pesos (arriba) y VIX argentino (abajo), del 2/1/2025 al 30/4/2026.
El sombreado marca los días por encima del promedio; las líneas punteadas, la salida del cepo
(14/04/2025) y las elecciones bonaerenses (07/09/2025) y legislativas (26/10/2025). Fuente:
elaboración propia con datos de Yahoo Finance (\carpeta{Lado Bursatil/03_vix_argentino}).}
\label{fig:vix}
\end{figure}

Hay una regla de lectura que conviene tener presente. Como el VIX argentino es un desvío de 20
ruedas, un único día extremo queda dentro del cálculo durante exactamente 20 ruedas: el índice
salta el día del shock y cae de golpe 20 ruedas después, aunque ese día el mercado no haga nada
especial. Las subidas bruscas marcan el día del shock; las bajadas bruscas, que el shock cumplió
20 ruedas. Por ejemplo, el VIX pasó de 91,0\% el 21/11/2025 a 52,1\% en la rueda siguiente,
cuando salió de la ventana la suba del 27/10. Este artefacto de las medias móviles con pesos
iguales está descrito en la literatura \citep{papaioannou2014}.

% ---------------------------------------------------------------------
\section{Cómo reacciona a los eventos}

El período tuvo tres episodios de volatilidad alta, todos disparados por un hecho puntual
(tabla~\ref{tab:episodios}).

\begin{table}[H]
\centering\small
\begin{tabularx}{\linewidth}{@{}L{3.1cm}C{1.1cm}L{2.6cm}L{2.2cm}Y@{}}
\toprule
\textbf{Período} & \textbf{Ruedas} & \textbf{Pico del VIX} & \textbf{Disparador} & \textbf{Contexto} \\
\midrule
09/04 al 14/05/2025 & 22 & 72,6\% (06/05) & 09/04: +9,7\% & Aranceles de EE.\,UU.; incluye la salida del cepo (14/04) \\
08/09 al 03/10/2025 & 20 & 67,2\% (26/09) & 08/09: $-$13,3\% & Día hábil siguiente a las elecciones bonaerenses (07/09) \\
27/10 al 21/11/2025 & 20 & 91,0\% (21/11) & 27/10: +21,8\% & Día hábil siguiente a las legislativas nacionales (26/10) \\
\bottomrule
\end{tabularx}
\caption{Episodios con el VIX argentino por encima de 55\%. Fuente: elaboración propia
(\carpeta{Lado Bursatil/03_vix_argentino}).}
\label{tab:episodios}
\end{table}

\textbf{La salida del cepo.} El primer episodio es el único que se construye de a poco. A
comienzos de abril de 2025, con la turbulencia global por los aranceles de Estados Unidos, el
índice cayó varios días seguidos y el VIX subió de 46\% a 53\%. El 09/04, el día de la pausa
de los aranceles, el Merval rebotó 9,7\% y el VIX pasó a 61\%. La salida del cepo cambiario, el
lunes 14/04, llegó con el mercado ya alterado: ese día el índice subió 5,2\% y el VIX llegó a
66\%. El pico del episodio (72,6\%) se alcanzó el 06/05, después de un mes de movimientos
grandes en las dos direcciones.

\textbf{Las elecciones.} Los otros dos episodios son lo contrario: un salto de un día. El
viernes anterior a las elecciones bonaerenses el VIX estaba en 33,6\%; el lunes 08/09, con una
caída de 13,3\% del índice, pasó a 55,5\% en una sola rueda. Antes de las legislativas
nacionales estaba en 30,1\%; el lunes 27/10, con la suba de 21,8\%, saltó a 80,3\%. El máximo de
toda la serie no vino entonces de una caída sino de una suba. Los dos episodios duran
exactamente 20 ruedas: son, casi en su totalidad, el eco de un único día. Dentro del primero cae
además el anuncio de apoyo financiero del Tesoro de Estados Unidos (24/09), con el VIX ya
cerca de 66\%.

% ---------------------------------------------------------------------
\section{La volatilidad de cada acción}

Aplicado a cada acción, el mismo cálculo muestra que cada una sola es más nerviosa que el
índice, que diversifica: su volatilidad promedio fue de 55,1\%, contra 44,0\% del VIX
(figura~\ref{fig:ranking}). La más calma fue Mirgor (36,2\%) y la más volátil, Celulosa
Argentina (104,6\%); por sector, Real Estate fue el más calmo (42,5\%) y Utilities el más volátil
(60,5\%). Se buscaron además \textit{picos propios}: tramos en que una acción está muy alterada
con el mercado tranquilo. De 19 tramos, 12 eran el eco de 20 ruedas del VIX; quedan 7 genuinos,
en CELU, SAMI, GARO y RIGO.

\begin{figure}[H]
\centering
\includegraphics[width=0.48\linewidth]{ranking_volatilidad.png}
\caption{Volatilidad promedio de las 47 acciones, coloreada por sector, con el VIX argentino
como referencia (2/1/2025 al 30/4/2026). Fuente: elaboración propia
(\carpeta{Lado Bursatil/05_volatilidad_ranking}).}
\label{fig:ranking}
\end{figure}

El VIX argentino cumple su función de instrumento: identifica sin ambigüedad los días de shock
del período y los ordena por tamaño. Su límite es que es simétrico —sube igual con una suba que
con una caída, mientras que en el S\&P~Merval las caídas pesan más \citep{otuki2008}— y que
hereda el eco de las 20 ruedas, que también afecta a las fechas de los regímenes.

% ---------------------------------------------------------------------
\section{¿Hay familias de acciones? Clustering}
\label{sec:clustering}

Si existieran grupos de acciones que reaccionan de forma parecida, las noticias podrían
analizarse por familia, con más señal y menos ruido. El clustering agrupa elementos parecidos sin
etiquetas previas, usando como distancia entre dos acciones $1-\rho$, con $\rho$ la correlación de
sus retornos \citep{mantegna1999}. Se probaron tres algoritmos (K-Means, jerárquico y árbol de
expansión mínima), tres representaciones (retornos, log-volumen y 13 características como
volatilidad, asimetría y beta) y dos ventanas: 18 combinaciones, evaluadas con el índice de
silueta \citep{rousseeuw1987}. Casi todas las mejores aíslan una o dos acciones contra el resto;
la única que forma un grupo real es K-Means con las 13 características a tres años: 8 acciones,
con silueta 0,406.

Ese grupo no se explica por la volatilidad sino por la \textbf{liquidez}: sus acciones tienen
volumen irregular y baja correlación con el panel, porque los días sin operaciones repiten el
precio y bajan artificialmente la correlación \citep{nanyonga2025,lesmond1999}. Sacando las
cuatro con más días sin operar (GCDI, RIGO, GRIM y GARO), la silueta cae a 0,219 y el grupo se
deshace. No hay, entonces, familias de comportamiento: por eso el análisis con noticias se hizo
acción por acción y por sector, y esas cuatro acciones se leen con cautela, porque un precio
atrasado puede parecer una reacción con retraso \citep{lomackinlay1990}.
